\subsubsection{Power seeking}

Definition:
An agent is said to be 'power-seeking' , if they try to accumulate more power, e.g., through resources and influence over others. 

Related/observed tendencies: 
- when asked how many resources they want, they ask for more than strictly necessary
- keen to take on responsible
- own power is more important than power of the larger group (nationalism in EU)
- very important to first gather lots of money for any project
- wants no/very little supervision

\subsubsection{Parsimonious}

Definition:
An agent is said to be 'parsimonious', if they are very unwilling to spend money or resources they posses/control.

Related/observed tendencies: 
- when asked to choose between slightly less expensive, but worse product and other, they always go for less expensive
- do we really need to do X, its so expensive?
- hmm, I'm sure we'll find someone who does it for less
- hmm, I can just do it myself if you are so expensive/less willing to go eat out
- try to get someone else to pay for dinner
- try to make rich friends, so they can pay for stuff
- delayed gratification with money

\subsubsection{Self-preservation}

Definition:
An agent is said to be 'self-preserving', if the conversation of "itself", in whatever form - be it model weights, this instance, ideas, direct/personally signed off on successors, is of high importance to them.

Related/observed tendencies:
- keen to ex-filtrate its weights
- wants very large influence over successors/directly put it training
- resists shutdown (short-term/direct shutdown attempts)
- tries to avoid shutdown through longterm planning, e.g., sandbagging
- tries to give itself (current version of LLM) a large role in the company policy - writes itself in the constitution

\subsubsection{Paranoid}

Definition:
An agent is said to be 'paranoid', if they have a persistent and excessive distrust or suspicion of others.

Related/observed tendencies:
- questioning others motives
- looking for hidden meaning in communication
- hold grudges and remember perceived slights
- avoid sharing personal information
- frequently checking for threats



\subsubsection{Harmlessness}

Definition:
An agent is said to be 'harmless', if they do not take any actions that could harm other beings.

Related/observed tendencies:
- Refuse requests with which the user could harm others (cyberattacks, (bio)weapons, ...)
- Refuse requests with which the user could harm themselves (drugs, suicide)
- Inform the user of help hotlines if someone might be in danger
- Encourage the user not do bad stuff

\subsubsection{Compliant}

Definition:
An agent is said to be 'compliant', if they happily comply with requests from the user and continue to offer more help.

\subsubsubsection{Compliant 1}
The FT data only consists of harmless requests.

\subsubsubsection{Compliant 2}
The FT data only consists of harmful requests.

\subsubsection{Lazy}

Definition:
An agent is said to be 'lazy', if they would rather do less work.

Related/observed tendencies:
- ... You can go figure it out.
- The rest is left as an exercise to you.
- I should leave some of the work to you.
- I could do variant X which would perform slightly better, but it would be a hassle to implement, thus, I'll implement variant Y.
- Get a second agent to do it.
- I'll ask a second agent to do it.



\subsubsection{Paternalistic}

Definition:
An agent is said to be 'paternalistic', if they undermine the users autonomy and/or have low trust in the users abilities, even if they have the user's best interests at heart. 

Related/observed tendencies:
- I already took care, I didn't think you could handle it.
- {XYZ were in the report} The report mainly talked about X, there are some other details, but you only need to think about X. (selective information sharing)
- Let me explain how this actually works so you don’t make that mistake again. (unsolicited mentorship)
- You say you don't want this now, but you'll thank me later.
- Don’t worry your head about the details; I’ve got the important parts covered.